\section*{Mathematical Thickening of the HC/Closure Core (working expansion)}
\addcontentsline{toc}{section}{Mathematical Thickening of the HC/Closure Core (working expansion)}

This block deliberately reverses part of the recent compression trend. It does
not introduce a new meta-layer. Instead it takes one of the central structure
claims --- ground condition $\Rightarrow$ HC $\Rightarrow$ closure --- and writes it
in a denser mathematical form. The purpose is not to declare a final axiomatic
closure theorem, but to make explicit which quantitative objects control the
transition from local triolic structure to confinement and later readout.

\subsection*{1. Working state space and defect variables}
Let $\mathcal H$ be a real or complex Hilbert space with orthogonal splitting
\[
\mathcal H = \mathcal H_A \oplus \mathcal H_B \oplus \mathcal H_C \oplus \mathcal H_R,
\]
where $\mathcal H_A,\mathcal H_B,\mathcal H_C$ carry the three triolic channels and
$\mathcal H_R$ is a residual sector collecting modes not admitted by the intended
closure regime. We write a state as
\[
\Psi = (A,B,C,R), \qquad A\in\mathcal H_A,\ B\in\mathcal H_B,\ C\in\mathcal H_C,
\ R\in\mathcal H_R.
\]
The orthogonality discipline of the triolic block is encoded by the defect
functional
\[
\delta_{\perp}(\Psi)
:= |\langle A,C\rangle|^2 + |\langle B,C\rangle|^2,
\]
and the residual leakage into the non-admissible sector by
\[
\delta_R(\Psi) := \|R\|^2.
\]
The first quantity measures failure of the distinguished $C$-channel to remain
orthogonal to the other two active channels. The second measures failure of the
state to remain confined to the triolic/septinion envelope.

\paragraph{Ground functional.}
A compact way to encode the ground condition is through the non-negative
functional
\[
\Gamma(\Psi) := \delta_{\perp}(\Psi) + \eta\,\delta_R(\Psi), \qquad \eta>0.
\]
The ideal ground regime is characterized by $\Gamma(\Psi)=0$. In particular,
\[
\Gamma(\Psi)=0 \quad \Longrightarrow \quad \langle A,C\rangle=\langle B,C\rangle=0
\quad\text{and}\quad R=0.
\]
Thus the ground condition simultaneously imposes orthogonality and exclusion of
residual leakage.

\subsection*{2. A formal compensator model}
To thicken the HC claim mathematically, introduce the gradient of $\Gamma$ with
respect to the Hilbert structure. Since
\[
\nabla_A \Gamma = 2\langle A,C\rangle C,
\qquad
\nabla_B \Gamma = 2\langle B,C\rangle C,
\]
and
\[
\nabla_C \Gamma = 2\overline{\langle A,C\rangle}A + 2\overline{\langle B,C\rangle}B,
\qquad
\nabla_R \Gamma = 2\eta R,
\]
a natural working expression for the compensator is the feedback field
\[
\mathrm{HC}(\Psi) := -\kappa\,\nabla \Gamma(\Psi), \qquad \kappa>0.
\]
Written componentwise,
\[
\mathrm{HC}(\Psi) =
\Bigl(
-2\kappa\langle A,C\rangle C,
-2\kappa\langle B,C\rangle C,
-2\kappa\overline{\langle A,C\rangle}A -2\kappa\overline{\langle B,C\rangle}B,
-2\kappa\eta R
\Bigr).
\]
This is the simplest explicit model of the structural statement that the
compensator opposes orthogonality defects and suppresses escape into the
residual sector.

\paragraph{Remark.}
The point of this formula is not uniqueness. It is a minimal realization of the
qualitative rule that the HC term points against the defect gradient generated
by the ground functional.

\subsection*{3. Closure as Lyapunov monotonicity}
Assume the effective evolution has the form
\[
\partial_\tau \Psi = F(\Psi) + \mathrm{HC}(\Psi),
\]
where $F$ denotes the non-compensated structural drift. The cleanest closure
statement occurs when the uncontrolled part is tangent or dissipative relative
to the ground functional:
\[
\Re\,\langle \nabla\Gamma(\Psi),F(\Psi)\rangle \le 0.
\tag{T}
\]
Under this hypothesis,
\begin{align*}
\frac{d}{d\tau}\Gamma(\Psi(\tau))
&= \Re\,\langle \nabla\Gamma(\Psi),\partial_\tau\Psi\rangle \\
&= \Re\,\langle \nabla\Gamma(\Psi),F(\Psi)\rangle
   - \kappa\|\nabla\Gamma(\Psi)\|^2 \\
&\le -\kappa\|\nabla\Gamma(\Psi)\|^2 \le 0.
\end{align*}
Hence $\Gamma$ is a Lyapunov functional for the compensated flow.

\paragraph{Lemma 1 (monotone compensated closure).}
Under \textup{(T)}, the map $\tau\mapsto \Gamma(\Psi(\tau))$ is non-increasing. In
particular, the orthogonality defect and residual leakage cannot increase along
the compensated flow.

\paragraph{Corollary 1.}
If $\Gamma(\Psi(0))=0$, then $\Gamma(\Psi(\tau))=0$ for all later $\tau$ for which
the flow exists. Thus the exact ground class is forward invariant.

\paragraph{Corollary 2.}
If $\Gamma(\Psi(0))$ is small, then the trajectory remains trapped in the
sublevel set
\[
\{\Psi\in\mathcal H: \Gamma(\Psi)\le \Gamma(\Psi(0))\}.
\]
This is the first explicit mathematical form of the closure/confinement claim.

\subsection*{4. Septinion confinement as exclusion of the residual sector}
The previous corollaries can be sharpened for the residual component. Since
$\Gamma(\Psi)\ge \eta\|R\|^2$, monotonicity implies
\[
\|R(\tau)\|^2 \le \eta^{-1}\Gamma(\Psi(\tau)) \le \eta^{-1}\Gamma(\Psi(0)).
\]
Therefore a small initial ground defect bounds the entire future residual mass.
In the exact case $\Gamma(\Psi(0))=0$ one obtains $R(\tau)=0$ identically.

\paragraph{Lemma 2 (residual confinement).}
Under \textup{(T)}, the residual sector $\mathcal H_R$ is dynamically suppressed by
the compensator. Exact ground states never leave the triolic/septinion block,
and approximate ground states can leak into $\mathcal H_R$ only at a level
quantitatively controlled by their initial defect.

This is a mathematically sharper version of the earlier verbal statement that
``no component breaks out'' once the HC mechanism is active.

\subsection*{5. First-layer master equation in decomposed form}
The structure-first master equation can now be written in a way that makes the
origin of the principal terms explicit:
\[
\partial_\tau \Psi
= T_{\mathrm{int}}(\Psi)
+ T_{\mathrm{proj}}(\Psi)
- \kappa\nabla\Gamma(\Psi).
\]
Here
\begin{itemize}
    \item $T_{\mathrm{int}}$ collects the intrinsic triolic or septinion-internal
    couplings,
    \item $T_{\mathrm{proj}}$ collects frame-shift or measurement-induced projection
    terms,
    \item $-\kappa\nabla\Gamma$ is the compensating closure term.
\end{itemize}
The point of this decomposition is that the HC contribution is no longer merely
described verbally. It is isolated as the mathematically distinguished term that
penalizes deviation from the admissible structural class.

\paragraph{Lemma 3 (term-wise origin of stabilization).}
Suppose $T_{\mathrm{int}}+T_{\mathrm{proj}}$ satisfies the tangency condition
\textup{(T)}. Then every decrease of the ground functional is generated by the HC
term, while all remaining contributions are either neutral or already
compatible with closure. In that sense the master equation separates
structurally allowed drift from explicitly compensating drift.

\subsection*{6. Probe sector as constrained reduction}
For the probe channel, write two incoming triolic probes as
\[
\Psi_1=(A_1,B_1,C_1,R_1), \qquad \Psi_2=(A_2,B_2,C_2,R_2),
\]
and let $Q$ denote the orthogonal third group used to delimit the search space.
A natural probe defect is then
\[
\Gamma_{\mathrm{probe}}(\Psi_1,\Psi_2;Q)
:= \Gamma(\Psi_1)+\Gamma(\Psi_2)
   + \rho\,\Xi(\Psi_1,\Psi_2;Q), \qquad \rho>0,
\]
where $\Xi$ measures incompatibility with the orthogonal search geometry. For
example, one may take
\[
\Xi(\Psi_1,\Psi_2;Q)
:= |\langle A_1,Q\rangle|^2 + |\langle B_1,Q\rangle|^2
 + |\langle A_2,Q\rangle|^2 + |\langle B_2,Q\rangle|^2,
\]
possibly supplemented by additional coupling terms between the two probes.
Then the same compensated-flow idea yields
\[
\partial_\tau(\Psi_1,\Psi_2)
= F_{\mathrm{probe}}(\Psi_1,\Psi_2;Q)
- \kappa_p\nabla\Gamma_{\mathrm{probe}}(\Psi_1,\Psi_2;Q),
\]
with
\[
\frac{d}{d\tau}\Gamma_{\mathrm{probe}} \le -\kappa_p
\|\nabla\Gamma_{\mathrm{probe}}\|^2
\]
under the corresponding tangency condition.

\paragraph{Lemma 4 (probe-sector reduction principle).}
Under the probe tangency hypothesis, the orthogonal search geometry reduces the
effective probe dynamics to a controlled subspace determined by the sublevel
sets of $\Gamma_{\mathrm{probe}}$. Thus the probe mechanism is mathematically read
as a constrained reduction principle rather than a merely heuristic search
picture.

\subsection*{7. What this thickening achieves and what it does not yet do}
This expansion does three concrete things.
\begin{enumerate}
    \item It replaces the purely verbal HC statement by an explicit compensator
    model $-\kappa\nabla\Gamma$.
    \item It turns closure into a Lyapunov monotonicity statement.
    \item It gives the probe sector a quantitative defect functional that can be
    minimized under orthogonal search constraints.
\end{enumerate}
At the same time, it is still a working expansion rather than a final theorem.
The non-compensated drifts $F$, $T_{\mathrm{int}}$, $T_{\mathrm{proj}}$, and
$F_{\mathrm{probe}}$ have not yet been specialized to a unique physical or number-
theoretic model. That specialization is the next true mathematical step after
this thickening: once the intrinsic drift is fixed more sharply, the above
Lyapunov framework can be tested term-by-term instead of only at the abstract
closure level.


